% id: 75-05
% book: 베이직쎈 확률과 통계 · 04 조건부확률
% section: 유형 042 사건의 독립과 종속
% figure: none
% answer: 
% answer_source: 
% variant_level: 0 (원본 전사)
% 이 파일은 build.mjs 가 items.json 에서 만든다. 고칠 때는 items.json 을 고친다.
\smash{\raisebox{1.5mm}{\dmkichul{베이직쎈 확률과 통계 75쪽 05번}}}
다음 중 두 사건 $A$, $B$에 대한 설명으로 옳은 것은? \dmcondfill{$0<\mathrm{P}(A)<1$, $0<\mathrm{P}(B)<1$}{}
\par\medskip\par\noindent\hangindent=1.4em\hangafter=1 {\small ①}\ $A$, $B$가 서로 독립이면 $\mathrm{P}(A\mid B)\ne\mathrm{P}(A\mid\comp{B})$이다.\par\noindent\hangindent=1.4em\hangafter=1 {\small ②}\ $A$, $B$가 서로 독립이면 $\mathrm{P}(A\mid B)=\mathrm{P}(A\mid A\cap B)$이다.\par\noindent\hangindent=1.4em\hangafter=1 {\small ③}\ $A$, $B$가 서로 독립이면 $\comp{A}$, $\comp{B}$는 서로 종속이다.\par\noindent\hangindent=1.4em\hangafter=1 {\small ④}\ $A$, $B$가 서로 배반사건이면 $A$, $B$는 서로 독립이다.\par\noindent\hangindent=1.4em\hangafter=1 {\small ⑤}\ $A$, $B$가 서로 배반사건이면 $\mathrm{P}(A\mid B)=\mathrm{P}(B\mid A)$이다.\par\medskip
